\begin{BedrockListedFormatDiagram}{SAVE/RESTORE Base Header}
\BedrockFormatRowRange{(:term:base.terms.header:)[63:0]}{63}{0}{%
\BedrockFormatReserved{(:term:base.terms.reserved:)}{12}
\BedrockFormatField{STATUS}{16}
\BedrockFormatField{FLAGS}{16}
\BedrockFormatReserved{(:term:base.terms.reserved:)}{9}
\BedrockFormatField{D}{1}
\BedrockFormatField{5}{1}
\BedrockFormatField{4}{1}
\BedrockFormatField{3}{1}
\BedrockFormatField{2}{1}
\BedrockFormatField{1}{1}
\BedrockFormatField{0}{1}
\BedrockFormatField{FMT}{4}
}
\BedrockFormatRowRange{bitmap[63:0]}{63}{0}{%
\BedrockFormatField{state block bitmap}{64}
}
\end{BedrockListedFormatDiagram}
{\small The labels \texttt{5} through \texttt{0} are the six \texttt{GS\_VALID} bits. SAVE writes
\texttt{GS\_VALID=0x3f} and \texttt{FMT=0}. The label \texttt{D} denotes \texttt{DFA}, which records the hidden
current double-fault-active bit.\par}
\begin{BedrockListedStructLayout}
  {SAVE/RESTORE Fixed Base Block}
  {8}
  {offset from the 4 KiB-aligned save-area base}
  {offsets increase downward}
\BedrockStructRow{0x000--0x00f}{%
  \BedrockStructHeaderField{4}{BASE HEADER\\\texttt{0x000} \textperiodcentered\ 8 bytes}
  \BedrockStructHeaderField{4}{STATE BLOCK BITMAP\\\texttt{0x008} \textperiodcentered\ 8 bytes}
}
\BedrockStructRow{0x010--0x04f}{%
  \BedrockStructSlotSeries{R}{0}{8}
}
\BedrockStructRow{0x050--0x08f}{%
  \BedrockStructSlotSeries{R}{8}{8}
}
\BedrockStructRow{0x090--0x0bf}{%
  \BedrockStructSlotSeries{GS}{0}{6}
  \BedrockStructAnnotationField{2}{RESTORE uses slots\\selected by \texttt{GS\_VALID}}
}
\BedrockStructTallRow{0x0c0+}{%
  \BedrockStructExtensionField{8}{EXTENSION STATE COMPONENTS\\CPUID-defined slots \textperiodcentered\ each component starts on a 64-byte boundary}
}
\BedrockStructFooter
  {Each \texttt{Rn} and \texttt{GSn} cell is one 8-byte slot.}
  {\texttt{Rn}=\texttt{0x010}+8n; \texttt{GSn}=\texttt{0x090}+8n.}
\end{BedrockListedStructLayout}
